\ifdefined\sourcedir\else
  \def\local{}
  \documentclass[9pt]{extbook}

\usepackage[margin=18mm]{geometry}
\usepackage{paracol}
\usepackage{microtype}
\usepackage[hidelinks, linktoc=all, colorlinks=true, linkcolor=blue]{hyperref}
\usepackage{needspace}
\usepackage{xcolor}
\usepackage{fancyhdr}
\usepackage{alphalph}
\usepackage{multicol}
\usepackage{tocloft}
\usepackage{xfootnotes}
%\let\xref\footnote
%\xfndebugtrue

\microtypesetup{nopatch=footnote}

\pagestyle{fancy}

\fancyhf{}                 % clear all header and footer fields
\fancyhead[LE,RO]{\thepage}   % page numbers outer edge
\fancyhead[LO]{\leftmark}     % left page, inner header
\fancyhead[RE]{\rightmark}    % right page, inner header

\fancyhead[LE]{}
\fancyhead[RO]{\rightmark}
\fancyhead[LO,RE]{\thepage}

\renewcommand{\thefootnote}{\alphalph{\value{footnote}}}

\setlength{\columnsep}{10mm}
\setlength{\parindent}{0pt}
\setlength{\parskip}{1pt}
\setlength{\belowfootnoteskip}{14pt}
% Width reserved for verse numbers (tune)

\newlength{\verselabelwidth}
\newlength{\verselabelgap}

% Measure width from the widest label you expect (Jeremiah: "52:34")
\settowidth{\verselabelwidth}{\scriptsize 52:34}

% Half-hanging label: half in margin, half in text
\newcommand{\verselabel}[1]{%
    \hspace*{-0.5\verselabelwidth}%
    {\textsuperscript{\scriptsize #1}}%
}
\newcommand{\emptyverselabel}[1]{%
  \if\relax\detokenize{#1}\relax
    \hspace*{-0.5\verselabelwidth}%
    {\textsuperscript{\scriptsize\textcolor{gray}{---}}}%
  \else
    \hspace*{-0.5\verselabelwidth}%
    {\textsuperscript{\scriptsize\textcolor{gray}{#1}}}%
  \fi
}
\newcommand{\smartverse}[2]{%
\noindent%
\if\relax\detokenize{#2}\relax
\emptyverselabel{#1}
\else
\verselabel{#1}
\fi
\,#2\par
}

\newcommand{\VersePair}[4]{%
  \smartverse{#1}{#2}
  \switchcolumn[1]
  \smartverse{#3}{#4}
  \switchcolumn[0]*
}
\newcommand{\VerseTriple}[6]{%
  \smartverse{#1}{#2}
  \switchcolumn[1]
  \smartverse{#3}{#4}
  \switchcolumn[2]
  \smartverse{#5}{#6}
  \switchcolumn[0]*
}

\newcommand{\ensurepar}{
\ifhmode\unskip\par\fi
}
\newcommand{\poemindent}[1]{\hspace*{#1em}}
% Start a new line iff we're currently mid-paragraph,
% then typeset one poetry line.
\newcommand{\poemline}[2]{%
  \ifnum#1>1 \penalty5000 \fi
  {\setlength{\parskip}{0pt}%
    \raggedright
    \hangindent=2em%
    \hangafter=1%
    \leftskip=#1em%
   \noindent#2\par}%
}

\newcommand{\BookHeading}[1]{%
  \cleardoublepage
  \chapter*{#1}
  \addcontentsline{toc}{chapter}{#1}
  \fancyhead[LE]{#1}
}
\newcommand{\ChapterHeading}[1]{%
  \vspace{1.5em}%
  \Needspace{7\baselineskip}%
  \begin{center}
    \large\bfseries Chapter #1\par
  \end{center}
  \vspace{0.75em}
  \fancyhead[RO]{Chapter #1}
}
\newcommand{\columnAHead}{%
    \hspace*{-0.5\verselabelwidth}%
    \small (WEB)\par
  \setcounter{footnote}{0}
%  \setcounter{xfootnote}{0}
}
\newcommand{\columnBHead}{%
    \hspace*{-0.5\verselabelwidth}%
    \small (LXX)\par
  \setcounter{footnote}{0}
%  \setcounter{xfootnote}{0}
}

\newcommand{\columnCHead}{%
    \hspace*{-0.5\verselabelwidth}%
    \small (AT)\par
}

\newcommand{\ColumnHeadings}[1][2]{
  \columnAHead\switchcolumn[1]\columnBHead
  \ifnum#1=3
  \switchcolumn[2]\columnCHead
  \fi
  \switchcolumn[0]*
}
\newcommand{\DescriptiveHeading}[1]{%
{\raggedright\par\vspace{-0.3\baselineskip}{\fontsize{9.5}{11.5}\selectfont\textbf{#1}}\par\vspace{0.6\baselineskip}}
}
\usepackage{fontspec}
\usepackage{polyglossia}
\setmainfont{Libertinus Serif} % missing character
\setmainfont{TeX Gyre Pagella} % best overall, I think. Good x height
   % makes it appear bigger and more readable
%\setmainfont{Junicode} % x height too small
%\setmainfont{EB Garamond} % x height too small; no sc variant

% tighten spacing a bit
\addfontfeatures{LetterSpace=-1}
\linespread{0.98}

\setmainlanguage{english}
\setotherlanguage{hebrew}
\newfontfamily\hebrewfont[
  BoldFont       = {Ezra SIL},
  ItalicFont     = {Ezra SIL},
  BoldItalicFont = {Ezra SIL}
  ]{Ezra SIL}
\setotherlanguage{greek}
\newfontfamily\greekfont[
  BoldFont       = {SBL Greek},
  ItalicFont     = {SBL Greek},
  BoldItalicFont = {SBL Greek}
]{SBL Greek}

\tolerance=1000
\emergencystretch=2em
\hbadness=5000

\widowpenalty=6000
\clubpenalty=3000

\newif\ifbritish

\ifdefined\usebritish
\britishtrue
\else
\britishfalse
\fi

\def\sourcedir{\ifbritish BE\else US\fi}

% Dot leaders for chapters
\renewcommand{\cftchapleader}{\cftdotfill{\cftdotsep}}
\renewcommand{\cftchapdotsep}{\cftdotsep}

\makeatletter
\newcommand{\TOCwithouttitle}{%
  \@starttoc{toc}%
}
\makeatother

\tolerance=1500
\emergencystretch=1em

\typeout{USEBRITISH DEFINED: \ifdefined\usebritish YES\else NO\fi}
\typeout{SOURCEDIR = [\sourcedir]}


  \begin{document}
  \def\propertitle{The Books of Poetry }
  \frontmatter
  \begin{titlepage}
\thispagestyle{empty}

\vspace*{\fill}

\begin{center}

{\Huge\bfseries \propertitle Parallel Edition\par}

\vspace{1.5em}

{\Large
Using text based on the World English Bible in parallel with a
slightly modernised version of Brenton's English translation of the
Septuagint\par}
\ifbritish (British English edition)\par\fi
\vspace{3em}

{\large Edited by\par}

\vspace{0.5em}

{\Large Glen Prideaux\par}

\end{center}

\vspace*{\fill}

\end{titlepage}
  \clearpage
\thispagestyle{empty}
\vspace*{\fill}

%\begin{center}
\small
\textcopyright\ 2026 Glen Prideaux

\vspace{1em}

This work is licensed under the\\
Creative Commons Attribution 4.0 International Licence.

\vspace{1em}

You are free to share and adapt this material,
provided appropriate credit is given.

\vspace{1em}

To view a copy of this licence, visit\\
\texttt{https://creativecommons.org/licenses/by/4.0/}

\vspace{2em}

\textbf{\propertitle Parallel Edition}

Using text based on Brenton's English translation of the Septuagint\\
in parallel with the World English Translation

\vspace{1em}

Edited by Glen Prideaux

\vspace{2em}

{
\setlength{\parindent}{1cm}%
\noindent Citation:\par
\indent Prideaux, Glen.\\
\indent \propertitle Parallel Edition: An English Alignment of the Septuagint and Masoretic Text Traditions.\\
\indent Version 1.0. Zenodo, 2026.\\
\indent 

\noindent Short citation form:\par
\indent Prideaux 2026, \propertitle Parallel Edition.
}

\vspace{2em}

Scripture texts reproduced in accordance with the respective licences
of their source translations.

\vspace{2em}

Typeset in \LaTeX{} using the \TeX{} Gyre Pagella typeface for English
text and Ezra SIL for Hebrew.

This edition is maintained as an open scholarly project.
Source files, version history, and development documentation are available at:

\href{https://github.com/GlenPrideaux/OT\_parallel}{https://github.com/GlenPrideaux/OT\_parallel}


%\end{center}
  % intro.tex
\cleardoublepage
\section*{Foreword}
{\large
{
\setlength{\parindent}{2em}
I was reading recently in Romans and noticed a footnote in my Bible that said, ``Romans 3:14 Psalm 10:7 (see Septuagint)''. It got me wondering, just how different \textit{is} the Septuagint? I've done a bit of work with the Greek translation of the Hebrew scriptures, but never really done a side by side comparison. Having recently done some other work that involved putting copies of a couple of individual books side by side, with a focus on where there are known significant differences, I decided to extend that and produce a parallel version of the whole of the Hebrew scriptures in English. This is the result.

The Greek version of the Hebrew scriptures tends to be overlooked in the Western Protestant tradition, but it is worth remembering that although the original text was undoubtedly written in the Hebrew language, the Septuagint was translated from these Hebrew scriptures hundreds of years before the MT was fully settled. Even though it is a translation it provides key insight into the history of the transmission of the text. Where there are differences between the MT and the LXX we should consider that the LXX may actually represent an earlier textual tradition. Indeed, the discovery of the Dead Sea Scrolls supports this with a number of documents in Hebrew that are closer to the LXX than they are to the modern MT.

We speak about ``the Septuagint'' as if it was monolithic, but in reality it is a collection of translations made over a period of time. The original seventy (or seventy-two) scholars were tasked with the translation of the Torah into Greek in Ptolemaic Egypt in the third century BC. Much of the story is legendary, such as it being completed by seventy two translators (six from each of the twelve tribes) in exactly 72 days, with the resulting translations being identical despite translators being isolated from one another. It is likely true that the Torah was translated in Alexandria: it reflects Egyptian Greek usage and was probably created for Greek-speaking Jews in Egypt as there were large numbers of Jews in the diaspora who no longer primarly spoke Hebrew. Importantly, this seventy scholars story applies only to the Torah or Pentatuch: Genesis, Exodus, Leviticus, Numbers and Deuteronomy. The rest of the Greek translation of the Hebrew scriptures was developed over the next centuries. The earliest recoverable Greek translation for most books is called the Old Greek (OG). For many books the OG is what people mean when they speak of ``the Septuagint'' (and is usually what is meant here, too). The character of the OG varies enormously: some books are very literal, others very free, some highly literary, others almost interlinear. For example, OG Isaiah is relatively free and interpretive; OG Proverbs is elegant and rhetorical; OG Job is shortened and literary; OG Daniel is extremely free and paraphrastic; and OG Ezekiel is often fairly literal. Then there are some Greek texts that systematically revise older Greek toward closer conformity with the Hebrew consonantal text. These date from as early as second century BC.

\subsection*{Sources}
My source text for a standard English Bible, translated almost exclusively from the Hebrew Masoretic Text (MT) that has been the standard source ever since Jerome translated the Latin Vulgate in the early 5th century, I have selected the World English Bible (WEB). This was selected because it has been placed in the public domain, meaning not only am I free to use it, but I am also free to distribute any derived work, like this one. The WEB comes in two flavours, with US or British English. This affects mainly spelling but there are a few other differences. It was simple enough for me to script using either version and as a result this document is also available in two versions. This is the
\ifbritish
British English
\else
US
\fi\ version. The WEB text has not been edited apart from some re-ordering necessary to present parallel passages together. Footnotes and poetry structure reproduce the notes and structure in the  WEB source text. A moderate number of cross references have been added using data from Glenton Jelbert, \url{https://glentonjelbert.com/introducing-the-cross-referenced-bible/}. There are much more comprehensive cross reference databases available but extensive cross referencing was not the purpose of this project.

\ifbritish
The WEB British English edition has frequent repeats of this footnote:
\begin{itemize}
\item\textbf{\small When rendered in ALL CAPITAL LETTERS, “LORD” or “GOD” is the translation of God’s Proper Name (Hebrew \texthebrew{ יהוה} usually pronounced Yahweh)}; or
\item\textbf{\small LORD or GOD in all caps is from the Hebrew \texthebrew{ יהוה} Yahweh except when otherwise noted as being from the short form \texthebrew{ יה} Yah}.
\end{itemize}
The sometimes frequent repetition of the identical footnote has been judged as unnecessary and these have been removed from this edition.
\fi

For the Septuagint I have used Brenton's 1851 English translation of the Septuagint. Brenton mostly used the Roman/Sixtine text published in 1587 under Pope Sixtus V, based chiefly on Codex Vaticanus but with many editorial decisions and later corrections. At the time, this was effectively the standard printed Septuagint. Brenton also consulted Holmes and Parsons, some variant readings, and occasionally other witnesses. But fundamentally Brenton translates the ecclesiastical printed LXX tradition available before modern textual criticism. 

This is a 19th century translation and showing its age, but its age also makes it in the public domain and so free to make derivative works and distribute. That was important to me because I am not up to making an entirely fresh English translation, nor do I want this work to be encumbered by copyright considerations that would prevent its free distribution. (For a more contemporary English translation of the Septuagint, check out the NETS project.) In order to make the text more accessible to a modern reader, I have scripted changes to replace the most archaic language with contemporary equivalents, and manually edited the result with reference to the Greek.

My reference for the original Greek is the 1935 Rahlfs LXX. Like the text Brenton used, Rahlfs is largely based on Codex Vaticanus (B), but it is a critical edition that also draws from the Codex Sinaiticus (S), the Codex Alexandrinus (A) and many secondary witnesses. His edition is eclectic, not simply reproducing one manuscript but choosing readings judged preferable. Brenton's text mostly matches the Rahlfs text but may deviate from it where there are different readings in different Greek manuscripts. This work is principally based on Brenton, and I have really only referred to the Greek where Brenton is unclear or uses words whose meaning has changed in the 175 years since Brenton published his work. As a result, this is something of a hybrid where some books are almost all based on Brenton with only an occasional word translated from Rahlfs, and other books having substantial passages updated.

Words in italics indicating translator's addition of words absent in the original but necessary for the English to make sense are mostly as in Brenton. I have made a small number of additions. I have made more deletions where rewording into more contemporary English has made the additional material unnecessary. Where I have made changes I have attempted to make these consistent with Brenton's approach. Many of the footnotes are similarly Brenton's, although I have added a fair number. Brenton's translation does contain errors; I have corrected these where I have spotted them, but my work is mostly to make the translation understandable to a modern reader and not a full revision of Brenton's work so it is likely that some of his errors remain, not to mention any errors that I may have introduced. Text set in small caps and descriptive headings are also as specified by Brenton.

With the exception of Jeremiah (see below) chapter divisions are as they appear in the WEB.

\subsection*{Comments on specific books}
Some books need special comment.

Chapters 36, 37, 38 and 39 of \textbf{Exodus} are quite different in the MT and LXX. I have reordered the LXX a little in an attempt to put parallel verses together. You can reconstruct the original order using the chapter and verse references attached to each verse.

The MT of the books of \textbf{1 and 2 Samuel} is particularly unreliable, and the Old Greek of the LXX has its own issues. I have set a third column in parallel with them that has my translation of the Palestinian text of 1 and 2 Samuel, based on an updated American Standard Version with adjustments made where the Palestinian text deviates from the MT. This column is labelled (P). I am indebted to E.H. Blackmore for the work he did on tabulating the differences between the MT and P texts.

The Greek of the book of \textbf{Esther} is substantially longer than the MT, with some extended additional sections. These additions are generally considered deuterocanonical and are usually separated out in a way that breaks the flow of the book. Esther is presented here with the additions in the order they appear in the Greek text. Additions are identified by a letter in place of a chapter number.

The Greek of the book of \textbf{Job} is particularly difficult and in places positively ungrammatical. Brenton attempted to smooth many of these difficulties; I have preferred to let them stand, and many of Brenton's additions have been removed to allow the reader to experience the Greek warts and all, so to speak.

The LXX version of \textbf{Jeremiah} is structured differently to the MT, with the oracles against the nations included in the middle where the MT has them at the end. This leaves us with the problem of what order to present the material here. Because most readers will have access to the order of the MT in their usual English Bible, the decision was made that in this edition Jeremiah would be presented in the order it appears in the LXX. Chapter and verse numbering follows
the Brenton LXX (including a/b suffixes where present). Verse matching between versions was prepared based on the information provided at \url{https://www.ccel.org/bible/brenton/Jeremiah/appendix.html}. 

The LXX version of \textbf{Daniel} is substantially longer than the MT.  The additions have been included in they order they are in the Greek. Unlike most of the Old Testament, where one Greek translation (the Septuagint) dominates, Daniel exists in two different Greek forms. The Old Greek version of Daniel, prepared 1st or 2nd century BC, roughly the same time as the rest of the LXX, was replaced with Theodotian's version around the second century AD. This is a revision of the Greek text to bring it more into conformity with the Hebrew and it does include the additional material that is absent from the Hebrew (and appears to have been originally composed in Greek). This is the version that Brenton translates, and hence is the version presented here as LXX. I present an AI-translated English version of the Old Greek as a third column, labelled (OG), with the translation based on a source obtained from \url{http://ccat.sas.upenn.edu/gopher/text/religion/biblical/lxxmorph/61.DanielOG.mlxx}. In addition to the Hebrew text, the Greek versions include three additions. Two can be considered separate works and are generally considered to be part of the Apocrypha: \textit{Susanna}\/ is sometimes included as Daniel 13, but sometimes comes first in the Greek text since it appears to be set earlier. \textit{Bel and the Dragon}\/ appears at the end of the book, often as Daniel 14. These both are really separate works. The third addition is a major insertion between 3:23 and 3:24, including the \textit{Prayer of Azariah} and the \textit{Song of the Three Holy Children.}\/ This was probably composed separately and later than the rest of Daniel but there is no evidence of it having been circulated as a separate work or works. Since my focus here is not to present the whole of the Apocrypha but to enable comparison between textual traditions of texts that appear in both, I have made the somewhat arbitrary decision not to include \textit{Susanna}\/ or \textit{Bel and the Dragon,}\/ but to leave the chapter 3 insertions in place.

\subsection*{Brenton footnote abbreviations}
As mentioned above, the footnotes to the LXX column are mostly unchanged from Brenton's translations, and any changes have followed his style. He uses the following abbreviations and symbols:\par
\begin{tabular}{ll}
  \textit{Heb.}&for Hebrew.\\
  \textit{Gr.}&for Greek. \\
  \textit{Lit.}&for Literally. \\
  \textit{q. d.}&for \textit{quasi dicat} (as if he were saying).\\
  \textit{Comp.}&for Compare. \\
  \textit{A. V.}&for Authorised Version. \\
  \textit{Alex.}&for Alexandrine Text. \\
  \textit{Ald.}&for Aldine Text.\\
  \textit{App.}&for Appendix. \\
  +&for Sign of addition. \\
  —&for Sign of ommission. \\
  \textit{sc.}&\textit{scilicet}, that is to say.
\end{tabular}\par
}

\subsection*{Ongoing issues}
For reasons currently obscure to the editor, footnotes occasionally (e.g. Isaiah 19:9) appear on the page after they are referenced.

\bigskip
\noindent Glen Prideaux\\
March 2026

\bigskip
\noindent\textbf{Left:} MT English (WEB)\\
\textbf{Right:} LXX English (Brenton, updated)\\
Where a third column exists it is as noted above and at the beginning of the book.

\cleardoublepage
\chapter*{Contents}
\begin{multicols}{2}
\TOCwithouttitle
\end{multicols}
}

  \mainmatter
\fi
{
\renewcommand{\bookintrotext}{%
\textit{
NOTE: Some of the Greek text of Job is convoluted and difficult to follow. Brenton's translation smooths out many of these difficulties by guessing at the meaning and inserting extra words to make the English clearer than the Greek text actually is. In this edition some of those interpolations have been removed and where the English text seems expecially cryptic and telegraphic it is because that is the nature of the underlying Greek text. 
}}

\input{\sourcedir/019-Job_parallel.tex}
}
{
\renewcommand{\literalchapter}{Psalm}
\renewcommand{\bookintrotext}{%
\textit{
NOTE: Many psalms have a heading that is not traditionally assigned a
verse number in English Bibles. In this edition these headings appear
as verse 0. It is also worth noting that the Greek translated in these
headings as ``Of David'' is in the dative case that would normally be
translated ``To David'' or perhaps ``For David'', but this is
transparently a rendition of the Hebrew \texthebrew{לְדָוִד} that, while also being
open to these alternative translations, reflects the normal Hebrew
possessive, ``Of David'' and so this traditional translation has been
retained.  } }
\input{\sourcedir/020-Psalms_parallel.tex}
}
\input{\sourcedir/021-Proverbs_parallel.tex}
\input{\sourcedir/022-Ecclesiastes_parallel.tex}
{
\renewcommand{\bookintrotext}{%
\textit{The Song of Songs is written as a poetic dialogue, but the original Hebrew text does not identify its speakers. The speaker headings in the World English Bible column are editorial aids and are not part of the biblical text itself.}
\par
\quad\textit{In many places, changes of speaker can be recognised in Hebrew through grammatical features such as gender, number, and forms of address, distinctions that are often less clear in English translation. These headings are intended to reflect those shifts and assist the reader in following the movement of the poem.}\par
\vspace{6pt}
\textbf{Note:}
\textit{The word rendered as “beloved” in the LXX column translates the Greek word \textgreek{ἀδελφιδός}, a diminutive derived from “brother”. It is used throughout the Song as a term of affection between lovers, but it is difficult to convey that nuance in English. Although the Hebrew uses its word for “beloved”, the Septuagint translators consistently render it with this distinctive Greek expression.}
}
\input{\sourcedir/023-Song_of_Solomon_parallel.tex} 
}
\ifdefined\local
  \end{document}
\fi
